Software energy consumption has become an increasingly important concern as
computing systems grow in scale and complexity. A growing ecosystem of
software-based power monitoring tools aims to make energy consumption
observable without requiring dedicated external measurement hardware.
However, these tools differ substantially in what they measure, how they
obtain or estimate energy consumption, how they attribute hardware-level
energy to software entities, and at what granularity and in which execution
environments they operate. These differences complicate both tool selection
and the interpretation and comparison of reported measurements.

This survey provides a systematic analysis of software power monitoring
tools and the principles underlying their operation. We first distinguish
hardware-level energy measurement and estimation from the subsequent
problem of attributing energy consumption to software activities. Based on
a systematic review of the literature and existing tools, we develop a
taxonomy covering measurement sources, energy and activity models,
attribution mechanisms, monitoring granularity, deployment environments,
and operational constraints. We use this taxonomy to comparatively analyze
representative tools and expose the assumptions and trade-offs underlying
different monitoring approaches. We further examine how software power
monitoring tools are validated, distinguishing agreement with system-level
energy measurements from the correctness of software-level energy
attribution. From this analysis, we derive practical guidance for selecting
monitoring approaches and identify open challenges concerning attribution
accuracy, measurement uncertainty, virtualization and multi-tenancy,
portability, monitoring overhead, and reproducible validation. The
resulting framework provides a common foundation for interpreting existing
tools and for designing and evaluating future software energy monitoring
approaches.
